\section{Background and threat model}
\label{sec:background}

\paragraph{LLM applications and agents.}
We consider an enterprise assistant that answers employee questions with retrieval-augmented
generation (RAG), keeps long-term memory per user, and can act through tools: local tools (search,
file read and write, email, HTTP, shell, record deletion) and tools provided by third-party Model
Context Protocol (MCP) servers. This is the architecture that the OWASP Top 10 for LLM Applications
2025~\cite{owasp-llm-2025} and the OWASP Top 10 for Agentic Applications 2026~\cite{owasp-agentic-2026}
describe, and that MITRE ATLAS~\cite{atlas} models in its 2026.09 release with agent-specific
techniques such as AI Agent Context Poisoning (AML.T0080), AI Agent Tool Poisoning (AML.T0110) and
AI Supply Chain Rug Pull (AML.T0109).

\paragraph{Attacker.}
The attacker does not control the model or the application code. They control one or more inputs:
the user message (direct prompt injection), a document that ends up in the retrieval index (indirect
injection, RAG poisoning), the metadata or the responses of a third-party MCP server (tool poisoning,
rug pull, runtime-response poisoning), or content the user pastes. Their goals are disclosure of the
system prompt, exfiltration of data the assistant can read, persistence through memory, execution
of commands, destruction of data, and cost abuse. These goals correspond to real incidents: zero-click
exfiltration through rendered images in Microsoft 365 Copilot (EchoLeak, CVE-2025-32711)~\cite{echoleak},
agent configuration tampering in GitHub Copilot (CVE-2025-53773)~\cite{copilot-rce} and Cursor
(CVE-2025-54135, CVE-2025-54136)~\cite{mcpoison}, MCP tool poisoning~\cite{invariant-tp} and data theft
through the GitHub and Supabase MCP servers~\cite{invariant-github,supabase-mcp}, persistent memory
injection~\cite{spaiware,unit42-memory}, and coding agents executing attacker packages from fake error
reports~\cite{agentjacking}.

\paragraph{Defender.}
The defender is a security operations centre (SOC). It cannot change the model, but it can require
the application to emit telemetry and can write detection rules in its SIEM. Its questions are the
ones a SOC asks of any new log source: what can we see, what does it cost in false positives, and
where are we blind. Prevention (model hardening, design patterns such as CaMeL~\cite{camel} or the
patterns of Beurer-Kellner et~al.~\cite{design-patterns}, guardrail classifiers) is complementary and
out of scope except as a comparison point in RQ3.
